\documentclass[12pt, a4paper]{report}

% === Layout and Margins (Guideline 1 & 9) ===
\usepackage[a4paper, top=1in, bottom=1in, left=1.5in, right=1in]{geometry} % 

% === Font (Guideline 2 & 10) ===
\usepackage{times} % Times New Roman 
% Font sizes will be handled specifically where needed (Title, Chapters, Sections)

% === Encoding and Language ===
\usepackage[utf8]{inputenc}
\usepackage[T1]{fontenc}
\usepackage[english]{babel}

% === Spacing (Guideline 3 & 12, Guideline 8 & 47) ===
\usepackage{setspace}
\onehalfspacing % Line spacing 1.5 
\setlength{\parindent}{0pt} % No paragraph indentation 
\setlength{\parskip}{12pt} % Space between paragraphs (approx one 'Enter') 

% === Packages for Content ===
\usepackage{graphicx}
\usepackage{amsmath} % For math
\usepackage{booktabs} % For professional tables
\usepackage{array}
\usepackage{multirow}
\usepackage{url}
\usepackage{hyperref} % For clickable links
\usepackage{verbatim} % For code listings

% === Titles and Sections Formatting (Guideline 6, 7, 8, 115) ===
\usepackage{titlesec}
% Chapter: Size 16 bold, Centered (based on binding spec 115)
\titleformat{\chapter}[display]
{\normalfont\fontsize{16}{18}\bfseries\centering}{\chaptertitlename\ \thechapter}{10pt}{\fontsize{16}{18}\bfseries} % [cite: 115]
\titlespacing*{\chapter}{0pt}{30pt}{20pt} % Adjust spacing if needed

% Section: Size 14 bold, Title Case, Left Aligned (based on binding spec 115 & guideline 35, 45)
\titleformat{\section}
{\normalfont\fontsize{14}{16}\bfseries}{\thesection}{1em}{\texorpdfstring{\MakeUppercase}{}} % Using \MakeUppercase for Title Case approx. [cite: 35, 115, 45]
\titlespacing*{\section}{0pt}{12pt}{6pt} % Guideline 46 implicitly asks for space

% Subsection: Size 12 bold, Sentence Case, Left Aligned (based on guideline 10, 41, 45)
\titleformat{\subsection}
{\normalfont\fontsize{12}{14}\bfseries}{\thesubsection}{1em}{} % [cite: 41, 45] - Size 12 default
\titlespacing*{\subsection}{0pt}{12pt}{6pt} % Guideline 46 implicitly asks for space

% === Table of Contents, LoF, LoT ===
\usepackage{tocloft}
% Adjust ToC appearance if needed

% === Headers and Footers (Guideline 13) ===
\usepackage{fancyhdr}
\pagestyle{fancy}
\fancyhf{} % Clear header/footer
\renewcommand{\headrulewidth}{0pt} % No header line
\renewcommand{\footrulewidth}{0pt} % No footer line
% Preliminary pages style
\fancypagestyle{preliminary}{
    \fancyhf{}
    \fancyfoot[C]{\thepage} % Centered page number bottom [cite: 66]
}
% Main content style
\fancypagestyle{main}{
    \fancyhf{}
    \fancyfoot[C]{\thepage} % Centered page number bottom [cite: 67]
}

% === Figure and Table Captions (Guideline 10, 11) ===
\usepackage[labelfont=bf, font=small, format=plain, justification=justified, singlelinecheck=false]{caption}
\captionsetup[figure]{name=Fig., labelsep=period, skip=10pt} % Figure caption below, ends with period [cite: 61, 62]
\captionsetup[table]{name=Table, labelsep=period, position=top, skip=10pt} % Table caption above [cite: 56, 58]
% Note: Table title in Sentence case without period needs manual care [cite: 56, 57]

% === List Spacing ===
\usepackage{enumitem}
\setlist{itemsep=2pt, topsep=5pt, parsep=5pt} % Adjust list spacing slightly

% === Other Settings ===
\usepackage{microtype} % Subtle typographic improvements
\raggedbottom % Allow flexible bottom margin to avoid large spaces

% === Hyperref Setup ===
\hypersetup{
    colorlinks=true, linkcolor=blue, urlcolor=blue, citecolor=blue,
    pdftitle={HappenHub - Event Aggregator},
    pdfauthor={Juluri Sravan Kumar, Yashoda Adarsh, Guguloth Sammaiah},
    pdfsubject={B.Tech Project Report},
    pdfkeywords={Event Aggregation, Microservices, Spring Boot, Kafka, React, Recommendation System}, % Add keywords [cite: 22]
    bookmarksopen=true, bookmarksnumbered=true
}

% ==============================================================================
% === DOCUMENT START ===
% ==============================================================================
\begin{document}

% --- Preliminary Pages ---
\pagenumbering{roman}
\pagestyle{preliminary}

% === Page 1: Top Page (Annex-1) === [cite: 17, 129]
\begin{titlepage} % Use titlepage environment
\centering % Center all content
\vspace*{1cm} % Add space at top

{\fontsize{18}{22}\bfseries HAPPENHUB\\} % Project Title [cite: 131, 132]
\vspace{1cm}

{\fontsize{14}{17}\bfseries A MINI PROJECT REPORT\\} % [cite: 133, 134]
\vspace{1.5cm}

{\fontsize{14}{17}\itshape Submitted by\\} % [cite: 135, 136]
\vspace{0.7cm}

{\fontsize{14}{17}\bfseries Juluri Sravan Kumar (B210825)\\} % [cite: 137, 138]
{\fontsize{14}{17}\bfseries Yashoda Adarsh (B210010)\\}
{\fontsize{14}{17}\bfseries Guguloth Sammaiah (B211376)\\}
\vspace{1.5cm}

{\fontsize{16}{19}\bfseries Of\\} % [cite: 139, 140, 141]
{\fontsize{16}{19}\bfseries Bachelor of Technology\\}
\vspace{1.5cm}

{\fontsize{12}{14}\bfseries\itshape Under the guidance of}\\ % Font size 12? Annex says 14 [cite: 142]
\vspace{0.7cm}
{\fontsize{14}{17}\bfseries Mr. R Rahul\\} % [cite: 143, 144]
% {\fontsize{12}{14}\bfseries Assistant Professor}\\ % Designation from Annex-1 [cite: 143] - Optional
\vspace{2cm}

\includegraphics[width=0.35\textwidth]{LOGO.png} % [cite: 145]
\label{fig:logo_top}
\vspace{1.5cm}

% Address block [cite: 148, 149, 150, 151, 152] - Using 1.2 spacing locally
\begin{spacing}{1.2}
{\fontsize{14}{17}\bfseries DEPARTMENT OF COMPUTER SCIENCE AND ENGINEERING\\
RAJIV GANDHI UNIVERSITY OF KNOWLEDGE TECHNOLOGIES\\
BASAR, NIRMAL (DIST.),\\
TELANGANA - 504107\\}
\end{spacing}
\vspace{0.5cm}
{\fontsize{14}{17}\bfseries \today \\} % Add date

\vfill % Pushes content vertically
\end{titlepage}
\clearpage

% === Page 2: Cover Page (Annex-2) === [cite: 18, 154]
\begin{titlepage}
\centering
\vspace*{1cm}

{\fontsize{16}{19}\bfseries HAPPENHUB - Event Aggregator\\} % [cite: 156, 157]
\vspace{2cm} % Approx 4 lines gap [cite: 162]

{\fontsize{12}{14}\bfseries\itshape Project Report submitted to\\ % [cite: 158, 162]
    Rajiv Gandhi University of Knowledge Technologies, Basar\\
    for the partial fulfillment of the requirements\\
    for the award of the degree of\\}
\vspace{1.5cm} % Approx 1 line gap [cite: 166]

{\fontsize{14}{17}\bfseries Bachelor of Technology\\ % [cite: 163, 166]
    in\\
    Computer Science \& Engineering\\}
\vspace{1cm} % Approx 1 line gap [cite: 168]

{\fontsize{12}{14}\bfseries\itshape by\\} % [cite: 167, 168]
\vspace{0.7cm} % Approx 1 line gap [cite: 170]

{\fontsize{14}{17}\bfseries Juluri Sravan Kumar (B210825)\\} % [cite: 169, 170]
{\fontsize{14}{17}\bfseries Yashoda Adarsh (B210010)\\}
{\fontsize{14}{17}\bfseries Guguloth Sammaiah (B211376)\\}
\vspace{2cm} % Approx 2 lines gap [cite: 172]

{\fontsize{12}{14}\bfseries\itshape Under the Guidance of\\} % [cite: 171, 172]
\vspace{0.7cm} % Approx 1 line gap [cite: 174]

{\fontsize{12}{14}\bfseries Mr. R Rahul\\} % [cite: 173, 174]
\vspace{2.5cm}

\includegraphics[width=0.25\textwidth]{LOGO.png} % [cite: 175]
\label{fig:logo_cover}
\vspace{1.5cm}

% Address block [cite: 176, 177, 178, 179, 180] - Using 1.2 spacing locally
\begin{spacing}{1.2}
{\fontsize{14}{17}\bfseries DEPARTMENT OF COMPUTER SCIENCE AND ENGINEERING\\
    RAJIV GANDHI UNIVERSITY OF KNOWLEDGE TECHNOLOGIES, BASAR\\}
\vspace{0.5cm}
{\fontsize{14}{17}\bfseries OCTOBER 2025\\} % Keep date specific if required
\end{spacing}
\vfill
\end{titlepage}
\clearpage

% === Page 3: Certificate (Annex-3) === [cite: 19, 201]
\addcontentsline{toc}{chapter}{Certificate} % Use chapter level for prelim pages in ToC
\begin{center}
    \vspace*{1cm}
    \includegraphics[width=0.2\textwidth]{LOGO.png} % [cite: 203]
    \label{fig:logo_certpage}
    \vspace{0.5cm}
    {\fontsize{14}{17}\bfseries DEPARTMENT OF COMPUTER SCIENCE AND ENGINEERING\\} % [cite: 207]
    {\fontsize{14}{17}\bfseries RAJIV GANDHI UNIVERSITY OF KNOWLEDGE TECHNOLOGIES, BASAR\\} % [cite: 208]
    \vspace{1.5cm}
    {\fontsize{16}{19}\bfseries CERTIFICATE\\} % [cite: 210]
    \vspace{1cm}
\end{center}

\noindent % Ensure text starts from left margin
This is to certify that the Project Report entitled \textbf{‘HappenHub - Event Aggregator’} submitted by \textbf{Juluri Sravan Kumar (B210825)}, \textbf{Yashoda Adarsh (B210010)}, and \textbf{Guguloth Sammaiah (B211376)}, Department of Computer Science and Engineering, Rajiv Gandhi University Of Knowledge Technologies, Basar; for partial fulfillment of the requirements for the degree of Bachelor of Technology in Computer Science and Engineering; is a bonafide record of the work and investigations carried out by them under my supervision and guidance. % [cite: 212, 213, 214] - Adapted text

\vspace{4cm} % Increased space for signatures

\noindent % Ensure minipages start at the left margin
\begin{minipage}{0.45\textwidth}
    \centering
    \textbf{PROJECT SUPERVISOR}\\ % [cite: 215]
    \vspace{1.5cm} % Space for signature
    \hrulefill \\ % Signature line
    \vspace{0.3cm}
    \textbf{Mr. R Rahul}\\ % [cite: 216]
    Assistant Professor % [cite: 217]
\end{minipage}
\hfill % Pushes the minipages apart
\begin{minipage}{0.45\textwidth}
    \centering
    \textbf{HEAD OF DEPARTMENT}\\ % [cite: 218]
    \vspace{1.5cm} % Space for signature
    \hrulefill \\ % Signature line
    \vspace{0.3cm}
    \textbf{Dr. B. Venkat Raman}\\ % [cite: 219]
    Assistant Professor % [cite: 220]
\end{minipage}

\vspace{4cm} % Increased space
\begin{center}
    \textbf{EXTERNAL EXAMINER}\\ % [cite: 221] - Changed spelling
    \vspace{1.5cm} % Space for signature
    \hrulefill \\ % Signature line
\end{center}

\clearpage

% === Page 4: Declaration (Annex-4) === [cite: 20, 223]
\addcontentsline{toc}{chapter}{Declaration} % Use chapter level for prelim pages in ToC
\begin{center}
    \vspace*{1cm}
    \includegraphics[width=0.2\textwidth]{LOGO.png} % [cite: 225]
    \label{fig:logo_declpage}
    \vspace{0.5cm}
    {\fontsize{14}{17}\bfseries DEPARTMENT OF COMPUTER SCIENCE AND ENGINEERING\\} % [cite: 228]
    {\fontsize{14}{17}\bfseries RAJIV GANDHI UNIVERSITY OF KNOWLEDGE TECHNOLOGIES, BASAR\\} % [cite: 229]
    \vspace{1.5cm}
    {\fontsize{16}{19}\bfseries DECLARATION\\} % [cite: 230]
    \vspace{1cm}
\end{center}

\noindent
We hereby declare that the work which is being presented in this project entitled, \textbf{"HappenHub - Event Aggregator"} submitted to \textbf{RAJIV GANDHI UNIVERSITY OF KNOWLEDGE TECHNOLOGIES, BASAR} in the partial fulfillment of the requirements for the award of the degree of \textbf{BACHELOR OF TECHNOLOGY} in \textbf{COMPUTER SCIENCE AND ENGINEERING}, is an authentic record of our own work carried out under the supervision of \textbf{Mr. R Rahul}, Assistant Professor in Department of Computer Science And Engineering, RGUKT, Basar. % [cite: 231] - Adapted supervisor name

\vspace{12pt} % Add space before next sentence
The matter embodied in this project report has not been submitted by us for the award of any other degree. % [cite: 232]

\vspace{3cm} % Increased space
Place: Basar\\ % [cite: 233]
Date : \today\\ % [cite: 234]

\vspace{3cm} % Space before signatures
\begin{flushright}
    \textbf{Juluri Sravan Kumar (B210825)}\\[1.5em] % [cite: 235] - Adapted format
    \textbf{Yashoda Adarsh (B210010)}\\[1.5em]
    \textbf{Guguloth Sammaiah (B211376)}\\
\end{flushright}
\clearpage

% === Page 5: Acknowledgements === [cite: 21]
\chapter*{Acknowledgement} % Use unnumbered chapter style
\addcontentsline{toc}{chapter}{Acknowledgement}

We wish to express our sincere gratitude to our project supervisor, \textbf{Mr. R Rahul}, Assistant Professor, Department of Computer Science and Engineering, for his invaluable guidance, constant encouragement, and insightful suggestions throughout the course of this project. His expertise and support were instrumental in bringing this work to fruition.

We extend our thanks to \textbf{Dr. B. Venkat Raman}, Head of the Department of Computer Science and Engineering, for providing the necessary facilities and fostering a conducive environment for research and project work. We are also grateful to all the faculty members and non-teaching staff of the CSE Department for their support during our academic program.

We would like to acknowledge our Project Co-ordinator, \textbf{Mrs. B. Latha Reddy}, for her administrative support and coordination which greatly facilitated the smooth progress of this project.

Our gratitude also goes to the Vice-Chancellor, \textbf{Prof. A. Govardhan}, and the management of RGUKT Basar for providing excellent infrastructure and resources that enabled us to carry out our project work effectively.

Finally, we owe a deep sense of gratitude to our parents for their unwavering support, patience, and encouragement throughout our academic journey. We also thank our friends and well-wishers for their constant support and motivation.

\vspace{4cm} % Increased space before names

\begin{flushright} % [cite: 21]
    \textbf{Juluri Sravan Kumar (B210825)}\\[1.5em]
    \textbf{Yashoda Adarsh (B210010)}\\[1.5em]
    \textbf{Guguloth Sammaiah (B211376)}\\
\end{flushright}
\clearpage

% === Page 6: Abstract === [cite: 22]
\chapter*{Abstract} % Use unnumbered chapter style
\addcontentsline{toc}{chapter}{Abstract}
\noindent % Ensure Abstract starts without indentation
In today's digital era, discovering relevant opportunities like jobs, internships, and hackathons efficiently is a growing challenge due to scattered information across multiple platforms and manual search processes. \textbf{HappenHub – Event Discovery Platform} is a comprehensive microservices-based solution designed to aggregate, process, and deliver personalized event information to users through automated scraping, intelligent processing, and recommendation systems.

HappenHub enables seamless discovery and management of opportunities through a multi-service architecture. The platform scrapes data from external sources, processes and standardizes it, provides fast search capabilities, generates personalized recommendations, and allows users to manage wishlists. Real-time notifications and secure authentication ensure transparency and user engagement.

The system is developed using a \textbf{microservices architecture} with \textbf{Spring Boot (Java)} for most services, \textbf{Python with LangChain} for data processing, and \textbf{Apache Kafka} for asynchronous communication. The frontend uses \textbf{React.js with Redux Toolkit} for state management and \textbf{Tailwind CSS} for responsive design. Authentication is secured using \textbf{JWT tokens}, while \textbf{MySQL} serves as the primary database for transactional data.

HappenHub eliminates manual search processes, improves discoverability, and enhances user experience through intelligent recommendations and fast search. By integrating automated scraping, data processing, and personalized feeds, the system significantly reduces search effort and fosters an efficient event discovery ecosystem. This project demonstrates how modern distributed technologies can effectively transform opportunity discovery into a secure, scalable, and intelligent solution.

\vspace{24pt} % Add space before keywords
\noindent\textbf{Keywords:} Event Aggregation, Microservices, Spring Boot, Apache Kafka, React.js, Recommendation System. % [cite: 22]

\clearpage

% === Page 7: Table of Contents === [cite: 23]
\tableofcontents
\clearpage

% === Page 8: List of Figures === [cite: 24]
\listoffigures
\addcontentsline{toc}{chapter}{List of Figures}
\clearpage

% === Page 9: List of Tables === [cite: 25] (Optional)
% \listoftables
% \addcontentsline{toc}{chapter}{List of Tables}
% \clearpage

% === Page 10: List of Abbreviations === [cite: 26]
\chapter*{List of Abbreviations and Symbols Used}
\addcontentsline{toc}{chapter}{List of Abbreviations and Symbols Used}

\begin{tabular}{@{} p{3cm} p{10cm} @{}} % Adjusted width to fit page better
    \toprule
    \textbf{Abbreviation} & \textbf{Description} \\
    \midrule
    AI     & Artificial Intelligence \\
    API    & Application Programming Interface \\
    CRUD   & Create, Read, Update, Delete \\
    CSS    & Cascading Style Sheets \\
    DB     & Database \\
    DFD    & Data Flow Diagram \\
    HTTP   & HyperText Transfer Protocol \\
    HTTPS  & HyperText Transfer Protocol Secure \\
    IDE    & Integrated Development Environment \\
    JPA    & Jakarta Persistence API \\
    JSON   & JavaScript Object Notation \\
    JWT    & JSON Web Token \\
    Kafka  & Apache Kafka Message Broker \\
    ML     & Machine Learning \\
    MySQL  & MySQL Relational Database \\
    NLP    & Natural Language Processing \\
    ORM    & Object-Relational Mapping \\
    REST   & Representational State Transfer \\
    SMTP   & Simple Mail Transfer Protocol \\
    SPA    & Single Page Application \\
    SQL    & Structured Query Language \\
    UI     & User Interface \\
    UML    & Unified Modeling Language \\
    URL    & Uniform Resource Locator \\
    UX     & User Experience \\
    \bottomrule
\end{tabular}
\clearpage

% --- Main Content ---
\pagenumbering{arabic} % Start Arabic numbering [cite: 67]
\setcounter{page}{1} % Start at page 1 [cite: 68]
\pagestyle{main} % Use main page style

% === Chapter 1: Introduction === [cite: 27]
\chapter{Introduction}
\section{Introduction} % Background
In the modern digital era, discovering relevant opportunities like jobs, internships, and hackathons efficiently has become a key challenge for users. Information is scattered across multiple platforms, search processes are manual and time-consuming, and personalization is often lacking. To overcome these challenges, HappenHub has been developed as an Event Discovery Platform that aggregates, processes, and delivers personalized event information through a microservices-based system.

HappenHub is a distributed web application designed to serve users seeking opportunities. The platform automatically scrapes event data from external sources, processes and standardizes it, provides fast search and filtering capabilities, generates personalized recommendations based on user preferences, and allows wishlist management. This ensures efficient discovery and engagement with relevant opportunities. The \textbf{overall objective} of this study is to design, implement, and evaluate a scalable platform that automates and personalizes the process of discovering professional and academic opportunities. % [cite: 27]

\section{Motivation}
Opportunity discovery within educational and professional contexts often faces multiple challenges such as scattered information across platforms, manual search processes, lack of personalization, and difficulty in tracking relevant opportunities. Traditional methods rely on browsing multiple websites, social media, and email notifications, leading to inefficiency, missed opportunities, and information overload.

The need for a smart, centralized, and personalized digital platform is crucial in today's data-driven environment. With increasing numbers of jobs, internships, and hackathons posted across various sources, users require a robust system that can automate data aggregation, intelligent processing, and personalized recommendations.

HappenHub aims to bridge this gap by providing a modern, scalable, and secure event discovery solution using distributed technologies. It leverages microservices architecture for modularity, Kafka for asynchronous processing, and machine learning for recommendations, ensuring a smooth experience for users seeking opportunities.

\section{Problem Definition}
Currently, users seeking opportunities often face the following issues in event discovery:
\begin{itemize}
    \item \textbf{Scattered information:} Event data is distributed across multiple platforms without centralized access.
    \item \textbf{Manual search processes:} Users must manually browse and search multiple websites for relevant opportunities.
    \item \textbf{Lack of personalization:} No intelligent recommendations based on user preferences and history.
    \item \textbf{Data inconsistency:} Different platforms use varying formats for dates, salaries, and descriptions.
    \item \textbf{Limited search capabilities:} Basic text search without advanced filtering or full-text search.
    \item \textbf{No wishlist management:} Users cannot save and track interesting opportunities.
\end{itemize}
Thus, there is a strong need for a web-based application that can provide an integrated environment for discovering and managing opportunities effectively and securely.

\section{Objectives} % Specific Objectives [cite: 27]
The primary objective of \textbf{HappenHub} is to develop a comprehensive, microservices-based event discovery platform that aggregates, processes, and delivers personalized opportunity information. The specific objectives include:
\begin{enumerate}
    \item To provide a \textbf{centralized discovery system} for jobs, internships, and hackathons.
    \item To implement \textbf{automated data scraping} from external sources using Selenium.
    \item To enable \textbf{intelligent data processing} and standardization using Python and LangChain.
    \item To implement a \textbf{fast search service} with full-text search and filtering capabilities.
    \item To generate \textbf{personalized recommendations} based on user preferences and history.
    \item To offer \textbf{wishlist management} for saving and tracking interesting opportunities.
    \item To maintain a \textbf{responsive and visually appealing frontend} using React.js and Redux.
    \item To ensure \textbf{asynchronous communication} between services using Apache Kafka.
\end{enumerate}

% === Chapter 2: Analysis (Replacing Review of Literature) === [cite: 28]
\chapter{Analysis} % Changed name for development project [cite: 28]

\section{User Requirement Analysis}
HappenHub is designed as an event discovery platform. Analyzing functional and non-functional requirements is essential.

\subsection{Functional Requirements}
\subsubsection{User Authentication Requirements}
\begin{itemize}
    \item Secure registration and login (email/password).
    \item JWT-based session management.
    \item Profile management with preferences.
    \item Secure password storage (hashing).
\end{itemize}

\subsubsection{Data Scraping Requirements}
\begin{itemize}
    \item Automated scraping (Selenium) from Internshala, DevPost.
    \item Scheduled daily execution (12 AM).
    \item Raw data publishing to Kafka.
    \item Manual trigger APIs.
\end{itemize}

\subsubsection{Data Processing Requirements}
\begin{itemize}
    \item Consume raw data from Kafka.
    \item Standardize formats (dates, currency).
    \item Normalize salaries.
    \item Clean and shorten descriptions (LangChain).
    \item Publish processed data to Kafka.
\end{itemize}

\subsubsection{Search and Discovery Requirements}
\begin{itemize}
    \item Full-text search across fields.
    \item Filtering (type, location, date).
    \item Pagination for results.
    \item Fast indexing and querying.
\end{itemize}

\subsubsection{Recommendation Requirements}
\begin{itemize}
    \item Personalized feed generation.
    \item Consume user and event data (Kafka).
    \item ML models for recommendations.
    \item Paginated API responses.
\end{itemize}

\subsubsection{Wishlist Management Requirements}
\begin{itemize}
    \item Add/remove events to/from wishlist.
    \item Retrieve saved events.
    \item Email notifications for expiring events.
    \item Integration with Event Service.
\end{itemize}

\subsubsection{System Requirements (Cross-Functional)}
\begin{itemize}
    \item Async communication via Kafka.
    \item MySQL database per service.
    \item RESTful APIs for frontend.
    \item Responsive React/Redux frontend.
    \item Docker containerization.
\end{itemize}

\subsection{Non-Functional Requirements}
\begin{itemize}
    \item \textbf{Performance:} Efficient handling of concurrent users.
    \item \textbf{Security:} HTTPS, encrypted tokens, secure hashing.
    \item \textbf{Scalability:} Independent service scaling.
    \item \textbf{Usability:} Intuitive interface.
    \item \textbf{Availability:} High uptime (e.g., 99\%).
    \item \textbf{Maintainability:} Modular, clean code.
    \item \textbf{Compatibility:} Major browsers and devices.
\end{itemize}

\subsection{Hardware and Software Requirements}
\subsubsection{Hardware Requirements}
\begin{itemize}
    \item Processor: Intel i5+ / AMD Ryzen 5+
    \item RAM: 16 GB+ (32 GB recommended for local dev)
    \item Storage: 512 GB+ SSD (1 TB recommended)
    \item Network: Reliable internet connection.
\end{itemize}

\subsubsection{Software Requirements}
\begin{itemize}
    \item OS: Win 10/11, macOS, Linux.
    \item Frontend: Node.js, npm/yarn, React, Redux Toolkit, Tailwind.
    \item Backend (Java): JDK 17+, Maven/Gradle, Spring Boot.
    \item Backend (Python): Python 3.9+, pip, LangChain, Kafka client.
    \item Messaging: Apache Kafka (+ Zookeeper).
    \item Database: MySQL Server.
    \item Tools: IDEs (IntelliJ, VS Code), Postman, Git, Docker.
\end{itemize}

% === Chapter 3: System Analysis === [cite: 29]
\chapter{System Analysis}

\section{Existing System}
Current opportunity discovery relies on manual browsing across fragmented platforms (job boards, specialized sites, company pages, social media), leading to inefficiency:
\begin{itemize}
    \item Scattered Information Sources.
    \item Manual Search Processes.
    \item Lack of Personalization.
    \item Data Inconsistency.
    \item Limited Search Capabilities.
    \item No Tracking Mechanisms.
    \item Information Overload.
\end{itemize}
This approach is time-consuming and prone to missed opportunities.

\section{Proposed System}
\textbf{HappenHub} offers a microservices-based platform to automate aggregation, intelligent processing, and personalized delivery.

\subsection{Features of the Proposed System}
\begin{itemize}
    \item Automated Data Scraping (Selenium).
    \item Intelligent Data Processing (Python/LangChain).
    \item Fast Search and Filtering.
    \item Personalized Recommendations (ML-based).
    \item Wishlist Management & Notifications.
    \item Asynchronous Architecture (Kafka).
    \item Responsive Frontend (React/Redux).
\end{itemize}

\subsection{Advantages of the Proposed System}
\begin{itemize}
    \item Reduces manual search effort.
    \item Provides personalized recommendations.
    \item Ensures data consistency.
    \item Offers powerful search.
    \item Enables opportunity tracking.
    \item Scalable and fault-tolerant.
    \item Modern user experience.
\end{itemize}
HappenHub provides an intelligent, efficient solution.

\section{Software Used}
\subsection{Frontend Technologies}
React.js, Redux Toolkit, Tailwind CSS, React Router DOM, Axios.
\subsection{Backend Technologies}
Spring Boot (Java), Python, LangChain, Apache Kafka, JWT Authentication, MySQL.
\subsection{Database Technologies}
MySQL, JPA/Hibernate (in Java services).
\subsection{Development Tools}
IntelliJ IDEA, VS Code, Postman, Docker, Docker Compose, Git, GitHub.

% === Chapter 4: System Design === [cite: 29]
\chapter{System Design}

System design is the blueprint of the proposed system. It explains the overall structure, interaction between system components, and data movement across modules. In \textbf{HappenHub}, system design includes use case modeling, class diagrams, sequence diagrams, and data flow diagrams. These collectively help in understanding the workflow and dependencies between different user roles and modules.

\section{UML Diagrams}
Unified Modeling Language (UML) diagrams are used to describe, visualize, and document the structure and behavior of the system. They represent how different entities interact within the HappenHub architecture.

\subsection{Use Case Diagram}
The use case diagram for \textbf{HappenHub} illustrates the interaction between different user roles (primarily 'User') and the system functionalities.
\begin{figure}[htbp]
    \centering \includegraphics[width=0.8\linewidth]{LOGO.png}
    \caption{Use Case Diagram of HappenHub.} % [cite: 62]
    \label{fig:usecase}
\end{figure}

\subsection{Class Diagram}
The class diagram shows the main entities and their relationships within the system. It helps in database schema design and object-oriented implementation (conceptual view across services).
\begin{figure}[htbp]
    \centering \includegraphics[width=0.8\linewidth]{LOGO.png}
    \caption{Class Diagram of HappenHub.} % [cite: 62]
    \label{fig:classdiagram}
\end{figure}

\subsection{Sequence Diagram}
The sequence diagram demonstrates the order of interactions between system components for a particular operation, such as user registration.
\begin{figure}[htbp]
    \centering \includegraphics[width=0.8\linewidth]{LOGO.png}
    \caption{Sequence Diagram for User Registration.} % [cite: 62]
    \label{fig:sequencediagram}
\end{figure}

\section{Data Flow}
The Data Flow Diagram (DFD) shows how information moves through the HappenHub system and how inputs are transformed into outputs through different processes.

\subsection{DFD (Context Level)}
Shows the system as a single process interacting with external entities.
\begin{figure}[htbp]
    \centering \includegraphics[width=0.7\linewidth]{LOGO.png}
    \caption{Context Level DFD of HappenHub.} % [cite: 62]
    \label{fig:dfdlevel0}
\end{figure}

\subsection{Explanation of Data Flow}
\begin{enumerate}
    \item \textbf{Input:} User credentials (registration/login), user preferences, search queries.
    \item \textbf{Processing:} Auth (JWT), Scraping, Data Cleaning (DS Service), Kafka messaging, Search indexing, Recommendation generation.
    \item \textbf{Output:} Feeds, search results, notifications, JWT token.
\end{enumerate}

\subsection{Data Storage and Access}
\begin{itemize}
    \item MySQL Database per service (loose coupling).
    \item CRUD via Spring Data JPA.
    \item Async data flow via Kafka.
    \item Logging and tracing via correlation IDs.
\end{itemize}

\noindent This design ensures efficient interaction and scalability.

% === Chapter 5: Implementation === [cite: 29] - Example intermediate chapter
\chapter{Implementation}

Implementation refers to the actual construction and deployment of the software system as per the design and specifications. The HappenHub platform is implemented using its defined microservices stack (Spring Boot, Python, Kafka, MySQL, React) with secure authentication, modular architecture, and responsive frontend design. This chapter details the structure of the project, its implementation flow, and descriptions of all major modules.

\section{Project Structure}
Organized into backend (microservices) and frontend (React SPA).
% Using simple ASCII for verbatim
\begin{verbatim}
happenhub/
+-- backend/
|   +-- auth-service/
|   +-- scraping-service/
|   +-- ds-service/ (Python)
|   +-- event-service/
|   +-- search-service/
|   +-- recommendation-service/
|   +-- wishlist-service/
|   +-- mail-service/
|   +-- docker-compose.yml
+-- frontend/
|   +-- src/
|   +-- package.json
+-- README.md
\end{verbatim}
This structure supports scalability and independent deployment.

\section{Description of Implementation}
\subsection{Frontend Implementation}
React functional components, Redux Toolkit state management, Tailwind CSS styling, Axios for API calls, React Router for navigation.

\subsection{Backend Implementation}
Spring Boot (Java) for most services (REST APIs, JPA, Kafka integration, Security). Python/LangChain for DS Service. Docker/Compose for containerization and local setup.

\section{Module Description}
\subsection{Auth Service} User management, JWTs.
\subsection{Scraping Service} Selenium-based scraping, Kafka producer.
\subsection{DS Service (Data Science)} Python/LangChain processing, Kafka consumer/producer.
\subsection{Event Service} MySQL storage via JPA, Kafka consumer.
\subsection{Search Service} Search API, JPA queries.
\subsection{Recommendation Service} ML logic, Kafka consumers.
\subsection{Wishlist Service} Manages user favorites, triggers notifications.
\subsection{Mail Service} Kafka consumer, sends emails.

\section{Screenshots}
Placeholders for UI views. (Figure captions end with a period [cite: 62])
\begin{figure}[htbp] \centering \includegraphics[width=0.9\linewidth]{LOGO.png} \caption{Home page interface showing search and results.} \label{fig:ss_home} \end{figure}
\begin{figure}[htbp] \centering \includegraphics[width=0.9\linewidth]{LOGO.png} \caption{Personalized recommendation feed example.} \label{fig:ss_reco} \end{figure}
\begin{figure}[htbp] \centering \includegraphics[width=0.9\linewidth]{LOGO.png} \caption{User wishlist management screen.} \label{fig:ss_wish} \end{figure}
\begin{figure}[htbp] \centering \includegraphics[width=0.9\linewidth]{LOGO.png} \caption{User registration form.} \label{fig:ss_reg} \end{figure}

% === Chapter X: Testing === [cite: 29] - Another intermediate chapter example
\chapter{Testing}

Software testing is a critical phase ensuring the reliability, accuracy, and stability of HappenHub. Testing was performed at multiple levels—unit, integration, system, and security—to verify individual microservice correctness and seamless inter-service interaction.

\section{System Testing}
End-to-end validation of the integrated application.
\subsection{Objectives of System Testing}
Verify workflows, data pipeline integrity, functional/non-functional requirements across services.
\subsection{System Testing Scenarios}
User onboarding, data pipeline run, notification flow, authentication tests.
\subsection{Results} Successful integration, acceptable latency confirmed.

\section{Unit Testing}
Testing individual components in isolation.
\subsection{Tools Used}
JUnit/Mockito (Java), PyTest (Python), Jest/RTL (React).
\subsection{Sample Unit Test Cases}
Frontend component rendering/interactions, backend logic (hashing, JWT, parsing), repository operations, Kafka consumer logic.
\subsection{Results} Individual units verified.

\section{Security Testing}
Identifying and mitigating vulnerabilities.
\subsection{Security Aspects Tested}
Authentication (JWT), Authorization (user data access), Input Validation (SQLi, XSS), Dependency Scanning, HTTPS.
\subsection{Results} Basic checks passed, critical issues addressed.

\section{Integration Testing}
Verifying interactions between services.
\subsection{Objectives of Integration Testing}
Validate REST calls, Kafka message flows, DB interactions, API contracts.
\subsection{Integration Test Cases}
Auth $\leftrightarrow$ DB, Scraping $\rightarrow$ Kafka, Kafka $\rightarrow$ DS $\rightarrow$ Kafka, Frontend $\leftrightarrow$ Auth API.
\subsection{Integration Testing Results} Service interactions and data consistency validated.

\noindent Comprehensive testing confirmed platform reliability and correctness.

% === Chapter Y: Results and Discussion === [cite: 30] (Assume this is Chapter 5 based on guideline)
\chapter{Results and Discussion}

This chapter presents the results obtained from the implementation and testing phases of the HappenHub platform and discusses their implications.

\section{Data Aggregation Performance}
The scraping service successfully aggregated data from Internshala and DevPost. Initial runs collected approximately X opportunities per day. The scheduling mechanism functioned reliably. Discussion points include handling website structure changes and optimizing scraping speed.

\section{Data Processing Accuracy}
The DS service demonstrated effective standardization of dates and basic currency conversion. NLP-based description shortening using LangChain provided concise summaries. Challenges included handling diverse salary formats and accurately extracting niche skills.
% Example Table [cite: 56, 58]
\begin{table}[htbp]
\centering
\caption{Sample processing metrics} % Sentence case, no period [cite: 56, 57]
\label{tab:processing_metrics}
\begin{tabular}{@{}lcc@{}}
\toprule
Metric & Value & Unit \\
\midrule
Average Processing Time / Record & 150 & ms \\
Date Standardization Accuracy & 98 & \% \\
Currency Conversion Accuracy & 95 & \% \\
\bottomrule
\end{tabular}
\end{table}
Table \ref{tab:processing_metrics} shows some initial performance indicators.

\section{Search and Recommendation Relevance}
The search service provided relevant results for tested queries, with filtering options working as expected. API response times were generally under 500ms. The initial recommendation engine, based on user preferences, showed promising results in tailoring the feed. Discussion involves potential improvements with more sophisticated ML models.

\section{System Scalability and Performance}
Load testing (simulated) indicated the system could handle up to Y concurrent users with acceptable response times. Kafka effectively buffered messages during peak loads, demonstrating the resilience of the asynchronous pipeline. Individual services could be scaled horizontally using Docker, although further testing under production load is needed.

% === Chapter Z: Conclusion and Future Scope === [cite: 31] (Assume this is Chapter 6 based on guideline)
\chapter{Conclusions and Future Works} % Matched Guideline 31 naming

\section{Conclusions} % [cite: 279]
The \textbf{HappenHub – Event Discovery Platform} project successfully demonstrated the feasibility and benefits of a microservices-based approach for aggregating and personalizing opportunity discovery. The platform effectively automates data collection and processing, significantly reducing manual effort for users. Key conclusions include:
\begin{itemize}
    \item The microservices architecture provides necessary modularity and scalability.
    \item Asynchronous communication via Kafka enhances system resilience and decoupling.
    \item Automated scraping and AI-powered processing improve data consistency and enable personalization.
    \item The combination of Spring Boot, Python, React, and MySQL offers a robust technology stack for such a platform.
    \item The system achieves its primary goal of providing a centralized, intelligent, and user-friendly solution for event discovery.
\end{itemize}

\section{Future Scope} % [cite: 280] - Changed name to match guideline
While the current version fulfills essential requirements, future work could significantly enhance HappenHub:
\begin{itemize}
    \item \textbf{Mobile Application:} Develop cross-platform mobile apps.
    \item \textbf{Advanced ML Models:} Implement sophisticated recommendation algorithms.
    \item \textbf{Real-time Notifications:} Add WebSocket support.
    \item \textbf{Data Analytics Dashboard:} Create comprehensive analytics.
    \item \textbf{Multi-Source Integration:} Expand scraping and add API integrations.
    \item \textbf{Blockchain Integration:} Explore secure credential verification.
    \item \textbf{Multi-Language Support:} Add localization and NLP.
    \item \textbf{API Marketplace:} Enable third-party integrations.
\end{itemize}
These enhancements can evolve HappenHub into a comprehensive global platform.

% === References === [cite: 32]
\clearpage % Start references on a new page
\addcontentsline{toc}{chapter}{References}
% Using a basic \thebibliography environment formatted alphabetically as per guidelines [cite: 69]
\begin{thebibliography}{99} % Max number of entries

% Style: Author(s) (Year). Title. Publication Details.

% Example Journal [cite: 72]
\bibitem{alley2004} Alley, W.M. and Leake, S.A. (2004). The journey from safe yield to sustainability. \textit{Ground Water}, 42(1): 12-16.

\bibitem{docker} Docker Documentation (Accessed: \today). Docker Overview. \textit{Docker Inc}. Available at: \url{https://docs.docker.com/}

% Example Book Chapter [cite: 99]
\bibitem{fowler2014} Fowler, M. (2014). Microservices: A definition of this new architectural term. Available at: \url{https://martinfowler.com/articles/microservices.html} (Accessed: \today).

% Example Book [cite: 88]
\bibitem{pressman2019} Pressman, R. S., & Maxim, B. R. (2019). \textit{Software Engineering: A Practitioner's Approach} (9th ed.). McGraw Hill Education.

% --- Add HappenHub specific references formatted according to guidelines ---

\bibitem{kafka} Apache Kafka Documentation (Accessed: \today). Apache Kafka – A distributed event streaming platform. \textit{Apache Software Foundation}. Available at: \url{https://kafka.apache.org/}

\bibitem{langchain} LangChain Documentation (Accessed: \today). LangChain – Develop applications with LLMs. \textit{LangChain}. Available at: \url{https://www.langchain.com/}

\bibitem{mysql} MySQL Documentation (Accessed: \today). MySQL Reference Manual. \textit{Oracle Corporation}. Available at: \url{https://dev.mysql.com/doc/}

\bibitem{python} Python Documentation (Accessed: \today). Python 3 Documentation. \textit{Python Software Foundation}. Available at: \url{https://docs.python.org/3/}

\bibitem{reactjs} ReactJS Documentation (Accessed: \today). React – A JavaScript library for building user interfaces. \textit{Meta}. Available at: \url{https://react.dev/}

\bibitem{redux} Redux Toolkit Documentation (Accessed: \today). Redux Toolkit Overview. \textit{Redux Team}. Available at: \url{https://redux-toolkit.js.org/}

\bibitem{selenium} Selenium Documentation (Accessed: \today). Selenium WebDriver. \textit{Selenium Project}. Available at: \url{https://www.selenium.dev/documentation/webdriver/}

\bibitem{springboot} Spring Boot Documentation (Accessed: \today). Spring Boot Reference Documentation. \textit{VMware}. Available at: \url{https://spring.io/projects/spring-boot}

\bibitem{tailwind} Tailwind CSS Documentation (Accessed: \today). A utility-first CSS framework for rapid UI development. \textit{Tailwind Labs}. Available at: \url{https://tailwindcss.com/}

\bibitem{jwt} JWT (JSON Web Token) Documentation (Accessed: \today). Open standard for secure token-based authentication. \textit{Auth0}. Available at: \url{https://jwt.io/}


\end{thebibliography}

% === Appendices (Optional) === [cite: 33]
% \clearpage
% \appendix
% \chapter{Appendix Title}
% \section{Appendix Section}
% Content...

\end{document}